\chapter*{阅读指引}

机器学习把经验转化为可以用于新输入的判断或行动。理解这一过程，不能只认识若干模型名称，
还要分清任务要求、数据提供的信息、模型能够表达的规律、学习算法怎样利用经验，以及最终用
什么证据评价结果。本部分建立这些贯穿全书的共同概念，暂不深入具体模型的计算细节。

《引言》从实际任务出发，说明机器学习系统由问题定义、数据、模型、学习过程和评价共同组成，
并区分训练、验证、测试与部署。《模型简介》承接其中的表示问题，比较不同模型怎样组织输入、
参数和预测关系；它帮助读者判断一个方法能表示什么，但不单独回答模型怎样获得知识。
《范式简介》转而考察经验的来源和使用方式，说明监督、自监督、强化学习及知识获取、演进与
迁移等路线怎样改变训练信息、反馈和评价条件。

三章共同给出后续内容的起点：任务规定要解决什么，模型规定用什么形式表达规律，学习范式
规定利用什么经验形成和更新规律。机器学习理论进一步研究在明确的抽样、损失和比较条件下，
能否证明算法用有限经验学到满足要求的结果；可信机器学习则把评价推进到实际使用，检查系统
能否适应应用中的人群、环境、权限和风险要求，并在条件变化或失效时得到发现与处置。
